\appendixsection{Appendix: Complements on linear representations}

The following proposition generalizes Prop. 13 of Chap. 1, §3, no. 8.

PROPOSITION 1. Let K be a commutative field, A a K-algebra, and V and W two left A-modules that are finite dimensional vector spaces over K. If there exists an extension L of K such that the \((A \otimes_{K} L)\) -modules \(V \otimes_{K} L\) and \(W \otimes_{K} L\) are isomorphic, then the A-modules V and W are isomorphic.

\begin{enumerate}
\def\labelenumi{\alph{enumi})}
\item
  Assume first that L is an extension of K of finite degree n.~Since \(V \otimes_{K} L\) and \(W \otimes_{K} L\) are isomorphic as \((A \otimes_{K} L)\) -modules, they are isomorphic as A-modules; but, as A-modules, they are isomorphic to \(V^{n}\) and \(W^{n}\) , respectively. Now V and W are A-modules of finite length; thus V (resp. W) is the direct sum of a family \((\mathrm{M}_{i}^{r_{i}})_{1 \leqslant i \leqslant p}\) (resp. \((\mathrm{N}_{j}^{s_{j}})_{1 \leqslant j \leqslant l}\) ) of submodules such that the \(M_{i}\) (resp. \(N_{j}\) ) are indecomposable, and that two of the \(M_{i}\) (resp. \(N_{j}\) ) with distinct indices are not isomorphic (Algebra, Chap. VIII, §2, no. 2, Th. 1). Then \(V^{n}\) (resp. \(W^{n}\) ) is the direct sum of the \(M_{i}^{nr_{i}}\) (resp. \(N_{j}^{ns_{j}}\) ); it follows (loc. cit.) that p = q and, after a suitable permutation of the \(N_{j}\) , that \(M_{i}\) is isomorphic to \(N_{i}\) and \(nr_{i}\) is equal to \(ns_{i}\) for \(1 \leqslant i \leqslant p\) . Thus, V is isomorphic to W.
\item
  Assume that K is an infinite field. The assumption implies that V and W have the same dimension over K. Let \((e_{i})_{1\leqslant i\leqslant m}\) and \((e_{i}^{\prime})_{1\leqslant i\leqslant m}\) be bases of V and W over K, and \((a_{\lambda})\) a basis of A over K. An isomorphism \(u:V\otimes_{K}L\to W\otimes_{K}L\) is a bijective L-linear map and at the same time an \((A\otimes_{K}L)\) -homomorphism, in other words it satisfies the conditions
\end{enumerate}

\[
a _ {\lambda} u \left(e _ {i}\right) = u \left(a _ {\lambda} e _ {i}\right) \quad \text { for   all } \lambda \text { and   all } i.\tag{1}
\]

Put \(a_{\lambda}e_{i} = \sum_{j}\gamma_{\lambda ij}e_{j}, a_{\lambda}e_{i}' = \sum_{j}\gamma_{\lambda ij}'e_{j}'\) , where the \(\gamma_{\lambda ij}\) and \(\gamma_{\lambda ij}'\) belong to K, and \(u(e_i) = \sum_{j}\xi_{ij}e_j'\) , where the \(\xi_{ij}\) belong to L. The conditions (1) can be written

\[
\sum_ {j} \xi_ {i j} \gamma_ {\lambda j k} ^ {\prime} = \sum_ {j} \gamma_ {\lambda i j} \xi_ {j k}\tag{2}
\]

for all \(\lambda, i, k\) . By assumption, the homogeneous linear equations (2) have a solution \((\xi_{ij}) \in L^{m^2}\) such that \(\det(\xi_{ij}) \neq 0\) . Since the coefficients of the system (2) belong to K, we know (Algebra, Chap. II, §8, no. 5, Prop. 6) that this system also admits non-trivial solutions in \(K^{m^2}\) ; let E be the vector subspace of \(\mathbf{M}_{m}(\mathrm{K})=\mathrm{K}^{m^{2}}\) , not reduced to zero, formed by these solutions. Let \((c_{l})_{1\leqslant l\leqslant p}\) be a basis of E, and put \((\xi_{ij})=\sum_{l}\eta_{l}c_{l}\) for any matrix \((\xi_{ij})\in\mathrm{E}\) ; then \(\det(\xi_{ij})\) is a polynomial \(\mathrm{P}(\eta_{1},\ldots,\eta_{p})\) with coefficients in K. On the other hand, we know (loc. cit.) that the solutions of (2) in \(L^{m^{2}}\) are of the form \(\sum_{l}\zeta_{l}c_{l}\) , this time with \(\zeta_{l}\in L\) ; for such a solution, \(\det(\xi_{ij})\) is equal to \(\mathrm{P}(\zeta_{1},\ldots,\zeta_{p})\) . Granting this, if we had \(\mathrm{P}(\eta_{1},\ldots,\eta_{p})=0\) for all \(\eta_{1},\ldots,\eta_{p}\in K\) , the coefficients of P would be zero since K is infinite; we would then have \(\mathrm{P}(\zeta_{1},\ldots,\zeta_{p})=0\) for all \(\zeta_{1},\ldots,\zeta_{p}\in L\) , which is contrary to our assumption. We can therefore find a matrix \((\xi_{ij})\in\mathrm{E}\) such that \(\det(\xi_{ij})\neq0\) , and the corresponding linear map \(V\to W\) is an isomorphism.

\begin{enumerate}
\def\labelenumi{\alph{enumi})}
\setcounter{enumi}{2}
\tightlist
\item
  General Case. Let \(\Omega\) be an algebraically closed extension of \(L\) , and \(K_0\) the algebraic closure of \(K\) in \(\Omega\) . The assumption implies that \(V \otimes_K \Omega\) and \(W \otimes_K \Omega\) are isomorphic ( \(A \otimes_K \Omega\) )-modules. Since \(K_0\) is infinite, part \(b\) ) shows that \(V \otimes_K K_0\) and \(W \otimes_K K_0\) are isomorphic ( \(A \otimes_K K_0\) )-modules. Retaining the notation in \(b\) ), the system (2) admits a solution ( \(\xi_{ij}\) ) \(\in K_0^{m^2}\) such that \(\det(\xi_{ij}) \neq 0\) . But the \(\xi_{ij}\) all belong to some algebraic extension \(K_1\) of finite degree over \(K\) . The ( \(A \otimes_K K_1\) )-modules \(V \otimes_K K_1\) and \(W \otimes_K K_1\) are isomorphic, and the proof is completed by using \(a\) ).
\end{enumerate}

PROPOSITION 2 (Maschke). Let A be a ring with unit element, E a left A-module, F a direct summand of E, G a group of finite order q, and \(\rho\) a linear representation of G on E. Assume that q.1 is invertible in A and that F is stable under G. Then there exists a complement of F in E that is stable under G.

Let p be a projection of E onto F. For all \(x \in E\) , put

\[
f (x) = q ^ {- 1} \sum_ {s \in G} \rho (s) ^ {- 1} p (\rho (s) x).
\]

We have \(f(x) \in \mathbf{F}\) and \(f(y) = y\) for all \(y \in \mathbf{F}\) , so \(f\) is a projection of \(\mathbf{E}\) onto \(\mathbf{F}\) . On the other hand, if \(t \in \mathbf{G}\) ,

\[
\begin{array}{r l} \rho (t) f (x) & = q ^ {- 1} \sum_ {s \in G} \rho (s t ^ {- 1}) ^ {- 1} p (\rho (s) x) \\ & = q ^ {- 1} \sum_ {s \in G} \rho (s) ^ {- 1} p (\rho (s t) x) \\ & = f (\rho (t) x). \end{array}
\]

Thus \(f\) commutes with \(\rho(\mathrm{G})\) , so \(\operatorname{Ker}(f)\) is a complement of \(\mathrm{F}\) stable under \(\mathrm{G}\) .

COROLLARY. Let G be a finite group of order q, and K a commutative field whose characteristic does not divide q. Then the group algebra of G over K is semi-simple.

Indeed, by Prop. 2, every module over this algebra is semi-simple.

PROPOSITION 3. Let A be a commutative ring, M an A-module, G a finite group acting on M, and A' an A-module. Assume that the order q of G is invertible in A. Let \(M^{G}\) be the set of elements of M invariant under G. Then the canonical homomorphism from \(M^{G} \otimes_{A} A'\) to \(M \otimes_{A} A'\) defines an isomorphism from \(M^{G} \otimes_{A} A'\) onto the module \((M \otimes_{A} A')^{G}\) of elements of \(M \otimes_{A} A'\) invariant under G.

Indeed, let \(\mathbf{Q}\) be the projection of \(\mathbf{M}\) onto \(\mathbf{M}^{\mathbf{G}}\) defined by \(\mathbf{Q}(x) = q^{-1}\sum_{g\in \mathbf{G}}g(x)\) for all \(x\in \mathbf{M}\) . If \(i\) denotes the canonical injection of \(\mathbf{M}^{\mathbf{G}}\) into \(\mathbf{M}\) , \(\mathbf{Q}\circ i\) is the identity map of \(\mathbf{M}^{\mathbf{G}}\) , so \((\mathbf{Q}\otimes 1_{\mathbf{A}'})\circ (i\otimes 1_{\mathbf{A}'})\) is the identity map of \(\mathbf{M}^{\mathbf{G}}\otimes_{\mathbf{A}}\mathbf{A}'\) . Since \(\mathbf{Q}\otimes 1_{\mathbf{A}'} = q^{-1}\sum_{g\in \mathbf{G}}(g\otimes 1_{\mathbf{A}'})\) , the image of \(i\otimes 1_{\mathbf{A}'}\) is \((\mathbf{M}\otimes \mathbf{A}')^{\mathbf{G}}\) . On the other hand, \(i\otimes 1_{\mathbf{A}'}\) is injective by what has been said before.

Remark. The preceding proposition applies in particular when \(A'\) is an \(A\) -algebra. In that case, \(M^G \otimes_A A'\) is an \(A'\) -submodule of \(M \otimes_A A'\) .

